\section{Evaluación y comparativas}

\subsection{Marco de evaluación en varias capas}

El presentador divide la evaluación en tres capas: comparativas fuera de línea (benchmarks), despliegue simulado (sandbox), repetición real (replay evaluation).

\begin{itemize}
    \item \textbf{Benchmark fuera de línea}: tareas del tipo D4RL, centradas en el normalized score.
    \item \textbf{Despliegue simulado}: usar un simulator para ejecutar la policy with hold-out reward to detect inflated Q-values.
    \item \textbf{Evaluación por repetición}: reescribir en el registro la trajectory producida por la policy y verificar si se cumplen reward/constraint.
\end{itemize}

\begin{importantbox}{Mientras más evaluación, mejor, pero debe dejarse espacio para human-in-the-loop}
El benchmark automático solo puede verificar capacity; para entender los failure modes también se necesita enlazar un human audit: observar el comportamiento de la policy en situaciones clave (edge case) y complementar con qualitative assessment.
\end{importantbox}

\begin{figure}[H]
\centering
\includegraphics[page=3,width=0.9\textwidth]{07_cs224r_offline_rl_2025.pdf}
\caption{Pipeline de evaluación en tres capas descrito en las Lecture slides\protect\footnotemark}
\end{figure}
\footnotetext{Slides página 3, resalta la escalera de evaluación de offline RL.}

\subsection{Tabla comparativa de la evaluación}

\begin{table}[H]
\centering
\begin{tabular}{p{0.2\textwidth} p{0.26\textwidth} p{0.26\textwidth}}
\toprule
Capa de evaluación & Propósito principal & Herramientas comunes \\
\midrule
Benchmark fuera de línea & Alinear rápidamente la capacidad global & D4RL, offline suites de OpenAI \\
Replay evaluation & Observar el reward de la policy dentro del replay buffer & replay simulator, human audit \\
Sandbox deployment & Usar un simulator para acercarse a un entorno cuasi productivo & simulator por dominio + twin models \\
\bottomrule
\end{tabular}
\caption{Tabla comparativa de los niveles de evaluación}
\end{table}

\subsection{Resumen de la sección}

La evaluación debe incluir tres capas: benchmark, intento en sandbox y verificación por repetición; cada capa puede capturar un modo de failure distinto. La revisión humana sigue siendo indispensable.

